\documentclass[a4paper]{article}
\usepackage{amsfonts}
\usepackage{amsmath}
\usepackage{amssymb}
\usepackage{graphicx}     

\begin{document}
  
\noindent \large Auteur: Hendrik De Bruyn \\
\noindent \large Studierichting: Informatica\\
\noindent \large Oefeningenreeks 3A nr. 16\\

\medskip

\normalsize

Beschouw de grafiek van de functie $f(x) = \dfrac{x^2(x+3)}{9}$. \\

\textbf{Opgave 16.1:} Bepaal de vergelijkingen van de raaklijnen in de snijpunten met de coördinaatassen. \\

\textbf{Oplossing 16.1:} \\

We beginnen met de afgeleide $f'(x)$ te berekenen. \\

$f(x) = \dfrac{x^2(x+3)}{9} = \dfrac{1}{9}x^3 + \dfrac{1}{3}x^2 \Rightarrow f'(x) = \dfrac{1}{3}x^2 + \dfrac{2}{3}x$ \\

De snijpunten met de X-as bevinden zich op $f(x) = 0$. We zien dan dat $x = 0$ of $x = -3$. \\

Het snijpunt met de Y-as bevind zich op $f(0)$. \\

De snijpunten zijn dus $(0, 0)$ en $(-3, 0)$. \\

Voor het snijpunt $(0, 0)$ berekenen we de richtingscoefficient als $f'(0) = 0$. \\

De raaklijn is dus $y - 0 = 0 \cdot (x - 0) \Rightarrow T_{(0, 0)} = 0$ \\

Voor het snijpunt $(-3, 0)$ berekenen we de richtingscoefficient als $f'(-3) = 1$. \\

De raaklijn is dus $y - 0 = 1 \cdot (x + 3) \Rightarrow T_{(-3, 0)} = x + 3$ \\

\textbf{Besluit 16.1:} De raaklijnen zijn dus $T_{(0, 0)} = 0$ en $T_{(-3, 0)} = x + 3$. \\

\newpage

\textbf{Opgave 16.2:} Bereken de coördinaten van de punten waar de raaklijn een hoek van $\dfrac{\pi}{4}$ met de X-as vormt. \\

\textbf{Oplossing 16.2:} \\

We berekenen de richtingscoefficient als $\tan \dfrac{\pi}{4} = 1 = f'(x)$. \\

We zoeken dus de $x$ waarvoor $f(x) = 1$. \\

$\Rightarrow \dfrac{1}{3}x^2 + \dfrac{2}{3}x - 1 = 0 \Rightarrow x^2 + 2x - 3 = 0 \Rightarrow x_{1,2} = \dfrac{-2 \pm 4}{2} = -1 \pm 2$ \\

Het eerste punt is $(-3, f(-3)) = (-3, 0)$. \\

Het tweede punt is $(1, f(1)) = \left(1, \dfrac{4}{9}\right)$. \\

\textbf{Besluit 16.2:} De coördinaten zijn dus $(-3, 0)$ en $\left(1, \dfrac{4}{9}\right)$. \\

\begin{thebibliography}{999}
%\bibitem{a} Referentie
\end{thebibliography}
\end{document} 
